%% v4/sec/intro.tex — §1 Introduction (TODAES long-form, post Round-19 pivot)
%%
%% Status (2026-05-16): Round-19 unlock — leader authorised §1 draft
%% (motivation + positioning sentence + C1-C6 contribution list).
%% §2 related-work three-way split is on hold pending dual-reviewer
%% closure (codex may challenge C1 first-analog claim).
%%
%% Target length: ~1.0 page in acmsmall single-column (~450-550 words).

\section{Introduction}
\label{sec:intro}

\subsection{Setting and gap}
\label{sec:intro-gap}

Large language models (LLMs) are now active participants in
electronic-design-automation (EDA) loops. In the digital and
high-level-synthesis (HLS) domains, recent ACM TODAES work has
established the template: an LLM proposes a design artefact, an EDA
tool returns ground-truth feedback (compile pass/fail, synthesis
metric, simulation trace), and the loop iterates until a contract
is met. AutoChip~\cite{autochip} drives this pattern for Verilog;
HLSRewriter~\cite{hlsrewriter}, C2HLSC~\cite{c2hlsc}, and
LHS~\cite{lhs} carry it through high-level-synthesis;
ConfiBench~\cite{confibench} adds principled confidence
benchmarking on top.

The \emph{analog} sizing loop is a natural next target: a SPICE
simulator already supplies the ground-truth oracle, and recent
work~\cite{analogcoder,adollm,eesizer,anaflow,ledro,autosizer,analogsage}
has shown that LLMs can iterate on netlists and revise device sizes.
Yet to our knowledge no prior work in a TODAES-class venue has
reported an analog-sizing LLM4EDA system evaluated on a real
foundry-bound process design kit (PDK). The combination of
\emph{industrial-node PDK confidentiality} and \emph{multi-vendor
frontier-LLM coverage} is the gap we close.

\subsection{Practical questions left unanswered}
\label{sec:intro-questions}

Three practical questions remain unanswered---the first two
constitute the headline conjunction; the third is the orthogonal
deployment-gate that any answer to (Q1)/(Q2) must clear.

\textbf{(Q1) Cost--quality choice.} Which LLM should a deployment
team pick under a token budget when running against a real
industrial PDK? Reported success rates on open benchmarks such as
AnalogGym~\cite{analoggym} cluster near saturation across many
checkpoints; the discriminating signal moves to dollars per run.
EEsizer~\cite{eesizer} took the first multi-LLM step but on
Ngspice-level open PDKs, with no cost--quality Pareto and no
per-LLM failure taxonomy.

\textbf{(Q2) Domestic-China viability.} Is domestic-China frontier
LLM inference (Moonshot, MiniMax, Xiaomi, DeepSeek) industrially
viable for analog sizing? To our knowledge zero published work
evaluates any of these four vendors on this workload, let alone all
four as a cohort.

\textbf{(Q3) PDK confidentiality (orthogonal gate).} Industrial
deployment is gated by an NDA that forbids foundry content from
leaving the controlled host. Concurrent work
AnalogAgent~\cite{analogagent} solves this by restricting inference
to small on-premise open-weight models (Qwen3 1.7B--14B). We take
the complementary path: admit any cloud LLM, prevent PDK content
from ever crossing the trust boundary. The two stories bracket the
industrial-deployment design space.

\subsection{This work}
\label{sec:intro-thiswork}

We present an eleven-LLM benchmark for industrial-node analog
sizing running on \textbf{Safe-Bridge}, a verifiable NDA-safe
wrapper around the PC\,$\leftrightarrow$\,remote-host SKILL IPC
channel. The eleven checkpoints span seven vendor families across
two countries: 6~US (Anthropic Opus~4.7 / Sonnet~4.6 / Haiku~4.5,
OpenAI GPT-5.5 / GPT-5.4-mini, Google Gemini~2.x) and 5~China
(Moonshot Kimi~K2.5, MiniMax~M2.7, Xiaomi MiMo~v2.5-pro, DeepSeek
V4-pro / V4-flash). The headline workload is LC\_VCO\_20\,GHz on a
foundry-bound PDK; an OpAmp transferability study (\S\ref{sec:methodology})
corroborates the platform on a second circuit topology. The
four-vendor China-frontier cohort turns domestic-China viability
from a single anecdote into a cross-vendor pattern.

\paragraph{Positioning.}
To our knowledge, this is the first NDA-safe, industrial-PDK,
eleven-LLM analog sizing benchmark in a TODAES-class venue,
extending prior LLM4EDA work (AutoChip~\cite{autochip},
HLSRewriter~\cite{hlsrewriter}, C2HLSC~\cite{c2hlsc},
LHS~\cite{lhs}, ConfiBench~\cite{confibench}) from the digital and
HLS domains into the analog / SPICE-loop domain under realistic
foundry-PDK confidentiality constraints. Concurrent TODAES analog
work AMSnet-KG~\cite{amsnetkg} targets open-PDK netlist
auto-design via knowledge-graph RAG; the two efforts occupy
orthogonal niches---dataset / retrieval methodology on open
substrate versus deployment substrate and cross-vendor evaluation
on closed substrate.

\paragraph{Contributions.}
\begin{itemize}
  \item \textbf{C1.}~First NDA-safe, closed-loop, industrial-PDK
    analog sizing benchmark comparing eleven LLM checkpoints under
    a scrubbed Safe-Bridge interface in ACM TODAES, extending the
    AutoChip / HLSRewriter / C2HLSC / LHS / ConfiBench digital-EDA
    precedents into the analog / SPICE-loop domain. Concurrent
    TODAES analog work AMSnet-KG~\cite{amsnetkg} targets open-PDK
    netlist auto-design via knowledge-graph RAG and is orthogonal
    along the NDA-PDK, scrub-substrate, and cross-vendor-cohort
    axes.
  \item \textbf{C2.}~\emph{Safe-Bridge} (\S\ref{sec:platform};
    overview in \S\ref{sec:overview}): a verifiable NDA-safe
    substrate for cloud-LLM-driven analog design on a real foundry
    PDK---code-anchored four-scrub-layer architecture with
    three-axis audit methodology, packaged in
    \S\ref{sec:reproducibility} as a reusable pattern for any
    LLM4EDA effort that must respect foundry NDA.
  \item \textbf{C3.}~An eleven-checkpoint industrial-node benchmark
    (\S\ref{sec:methodology}) spanning seven vendor families across
    two countries: the only cross-vendor cohort to evaluate frontier
    domestic-China LLM inference (Moonshot, MiniMax, Xiaomi,
    DeepSeek) on industrial analog sizing.
  \item \textbf{C4.}~A cost--quality Pareto with a labelled
    \emph{reasoning-token premium} sub-section
    (\S\ref{sec:results})---the sharpest single quantitative
    finding, characterising practitioner under-budgeting by
    $r/(1{-}r)$ for vendors that bill hidden reasoning tokens.
  \item \textbf{C5.}~An eight-bucket failure-mode taxonomy with
    per-LLM cross-vendor data (\S\ref{sec:failures}), moving
    beyond pass/fail headlines to characterise \emph{how} each
    vendor's behaviour breaks down at the margin.
  \item \textbf{C6.}~A reproducibility / artifact contribution
    (\S\ref{sec:reproducibility}): full software artifact
    structure, scrub-rule format, and replication recipe enabling
    future analog LLM4EDA evaluations on closed industrial PDKs.
\end{itemize}

\paragraph{Roadmap.}
\S\ref{sec:related} surveys related work across three clusters
(LLM for digital EDA, classical analog optimization, PDK-safe LLM
tooling). \S\ref{sec:overview} gives a one-figure system overview
of the agent loop, Safe-Bridge, and the spec-contract framework.
\S\ref{sec:platform} treats Safe-Bridge in depth: threat model,
four scrub layers, hard invariants, three-axis audit methodology.
\S\ref{sec:methodology} describes the benchmark methodology, the
eleven-checkpoint matrix, and threats to validity.
\S\ref{sec:results} reports the cost--quality Pareto and the
reasoning-token premium. \S\ref{sec:failures} presents the failure
taxonomy. \S\ref{sec:discussion} discusses cross-vendor
domestic-China viability. \S\ref{sec:reproducibility} packages the
artefact and replication recipe. \S\ref{sec:conclusion} concludes
and lists future work.
